\section{Methodology}
\label{sec:method}

\subsection{Problem formulation}
\label{sec:method:formulation}

We formulate cross-language scene text replacement in video as a
structured video-to-video transformation task. Given an input
video
\begin{equation}
  V = \{F_1, F_2, \dots, F_N\}, \quad
  F_t \in \mathbb{R}^{H \times W \times 3},
\end{equation}
we produce an output sequence
\begin{equation}
  V' = \{F'_1, F'_2, \dots, F'_N\},
\end{equation}
in which all scene text has been translated into a target language
while preserving geometric consistency, temporal coherence, and
photorealistic appearance.

We model each distinct text instance in the video as a
\emph{track}:
\begin{equation}
  T_i = \{Q_i^t, s_i, s'_i\}_{t=1}^{N},
  \label{eq:track}
\end{equation}
where $Q_i^t$ is the quadrilateral region of instance $i$ at
frame $t$ (undefined for frames in which the instance is not
visible), $s_i$ is the detected source-language string for that
instance, and $s'_i$ is its translation into the target language.
The pipeline described in the following subsections takes
$\{T_i\}$ as its central data structure; each stage enriches every
track with additional fields---homographies, edited ROIs,
propagated pixels, alpha masks---until the final compositing stage
produces $V'$.

\subsection{Pipeline overview}
\label{sec:method:overview}

Our system processes an input video through five sequential stages
(Figure~\ref{fig:pipeline}). Stage~1 (S1) detects scene text on a
regular sampling of frames, groups detections into temporally
coherent tracks via intersection-over-union matching, translates
each track's source text into the target language, selects one
frame per track as the \emph{reference} for that track, and fills
the remaining frames by optical-flow propagation of the detected
quad. Stage~2 (S2) computes a homography from each frame's quad to a
canonical frontal rectangle shared by the whole track, and refines
the result with a learned residual-homography predictor that
corrects tracker drift. Stage~3 (S3) warps the reference frame into
canonical space and passes it through a text-editing model that
accepts the target string and returns an edited region of interest
(ROI) in which the source text has been replaced. Stage~4 (S4)
adapts that single edited ROI to every other frame via a per-pixel
lighting correction and a per-frame blur prediction. Stage~5 (S5)
warps each adapted ROI back to the frame's perspective via the
inverse homography and composites it into the original frame.

\figplaceholderwide{pipeline}{\linewidth}%
  {Five-stage pipeline for cross-language scene text replacement.
   S1 detects, tracks, translates, and selects one reference frame
   per track; S2 computes and refines a homography to canonical
   frontal space; S3 runs the reference frame's canonical ROI
   through a text-editing model; S4 adapts the edited ROI to every
   other frame via a lighting correction module and our Blur
   Prediction Network; S5 warps each adapted ROI back to the
   original perspective and composites it into the frame.}%
  {fig:pipeline}

The system is organized around a single central dataclass,
\texttt{TextTrack}, that carries all per-track state through the
pipeline; each stage enriches this object rather than producing
parallel data structures. All tunable parameters are exposed via
YAML configuration files with validators that enforce weight-sum
constraints and disallow mutually incompatible combinations.
Expensive resources---the OCR engine, the translator, the
text-editing server, inpainters, and the two trained networks---are
lazily initialized on first use so that the pipeline starts quickly
and degrades gracefully when an optional component is missing.

The two models we trained---a Blur Prediction Network (BPN) in S4
and an Alignment Refiner in S2---are described in dedicated
subsections (\ref{sec:method:bpn} and~\ref{sec:method:refiner})
below. The remaining five ML models (PaddleOCR, CoTracker3,
AnyText2, SRNet, and Hi-SAM) are pretrained and used without
modification.

\subsection{Stage 1: Detection, tracking, translation, and reference selection}
\label{sec:method:s1}

S1 produces a set of \emph{tracks}, each a maximal temporal chain
of text detections that plausibly correspond to the same physical
text in the scene. We run PaddleOCR~\citep{paddleocr2024} on every
$N$-th frame (configurable) to obtain text quads, recognized source
strings, and per-detection confidences, and discard detections
whose recognized string has a Zipf frequency below a threshold
under \texttt{wordfreq}~\citep{wordfreq}---a simple but effective
filter against OCR false positives on textured or geometric
backgrounds. Frame-level detections are grouped into tracks
greedily by IoU match with a configurable break threshold; a
text-similarity check prevents accidental merging when two
unrelated quads briefly overlap.

For each grouped track we translate the source string once using
\texttt{deep-translator} with a Google translation backend and a
MyMemory fallback (with locale-code mapping, since MyMemory
expects \texttt{en-GB} rather than \texttt{en} and \texttt{zh-CN}
cased exactly so). Between OCR-detected frames, quads are filled by
optical flow. We support classical Farneback flow as a CPU
fallback, and CoTracker3~\citep{karaev2024cotracker3} online mode
as the default; the latter tracks the four corners of each quad
via a learned point tracker in sliding 60-frame windows, producing
substantially smoother trajectories than classical methods.

The final step of S1 selects, for each track, a \emph{reference
frame}: the frame at which we will perform the actual text edit.
Inspired by STRIVE~\citep{subramanian2021strive}, we apply a
four-filter funnel: (i)~keep detections with the maximum
recognized-text length within the track (to avoid partial
detections), (ii)~keep those with OCR confidence above a threshold,
(iii)~keep the top~$K$ by sharpness (measured as Laplacian
variance), and (iv)~rank the survivors by a weighted composite of
contrast (Otsu inter-class variance) and frontality (the ratio of
quad area to bounding-box area). This choice is critical: all of
S3's editing quality depends on the quality of the frame we hand
to the text-editing model.

\subsection{Stage 2: Frontalization and Alignment Refinement}
\label{sec:method:s2}

Given a track with quads in every frame and a single reference
frame, S2 establishes a \emph{canonical frontal space}: an
axis-aligned rectangle whose dimensions are derived from the
reference quad's edge lengths. For each detection we compute a
homography $H_{\mathrm{to\_frontal}}$ from the frame quad to this
canonical rectangle and store its inverse
$H_{\mathrm{from\_frontal}}$ on the same detection. S2 is
deliberately pure geometry: no pixels are warped here; warping is
done on demand by downstream stages. The baseline homography is
estimated via RANSAC with a 5-pixel reprojection threshold,
optionally using a grid of CoTracker points inside the quad for
multi-point least-squares refinement.

CoTracker's residual tracking error, while small (typically
4--8~pixels), causes pixel-level misalignment between the canonical
representation of one frame and another. This misalignment degrades
every downstream stage: lighting correction in S4 relies on
comparing pixel-aligned background regions, and S5's composite
inherits any drift introduced upstream. To address this we
designed and trained an \emph{Alignment Refiner}, a small
convolutional network that takes a reference canonical ROI and a
target canonical ROI and predicts four corner offsets. These
offsets are converted into a residual homography $\Delta H$ via a
differentiable direct linear transform (DLT), and folded into the
warp chain as $H_{\mathrm{to\_frontal}}' = \Delta H^{-1}\,
H_{\mathrm{to\_frontal}}$. The architecture and training scheme are
described in Section~\ref{sec:method:refiner}. A sanity gate
rejects predictions whose maximum corner offset exceeds 16~pixels,
silently falling back to the unrefined baseline in such cases.

\subsection{Stage 3: Text Editing}
\label{sec:method:s3}

S3 is the stage at which the target-language string actually
appears. Using $H_{\mathrm{to\_frontal}}$ from the refined S2
output, we warp the reference frame into canonical space to produce
a clean, rectified ROI. We optionally expand the ROI by 30\% per
side (padding with surrounding scene pixels) before editing, to
give the editor additional context for style matching, and then
crop the result back after editing.

The editor itself is accessed through a \texttt{BaseTextEditor}
abstract base class, allowing us to switch backends via
configuration. In our primary configuration the backend is
AnyText2~\citep{tuo2025anytext2}, a diffusion-based multilingual
text-image editor built on Stable
Diffusion~\citep{rombach2022stablediffusion}. We run the AnyText2
model in a separate Gradio server process and interact with it over
HTTP; this decouples the main pipeline from the editor's GPU memory
requirements and its own inference stack.

A practical issue arises when target and source differ substantially
in length. When, for example, a long source string is replaced by a
much shorter target, the source mask carries more space than the
target glyphs need, and AnyText2---like most mask-based scene text
editors---tends to fill the excess with hallucinated gibberish. We
address this with an \emph{adaptive mask} path. We estimate the
target string's natural width from a character-class width
heuristic: full width $1.0$ for CJK characters, $0.60$ for Latin
uppercase, $0.50$ for lowercase, $0.55$ for digits, and $0.30$ for
spaces, each multiplied by the canonical height. When this estimate
is more than $15\%$ narrower than the source canonical, we
(i)~pre-inpaint the source text out of the canonical ROI using the
pipeline's configured inpainter backend
(Section~\ref{sec:method:inpainter}); (ii)~restore a centered,
feathered strip of original pixels matching the target's natural
width; and (iii)~send this hybrid image, together with a shrunk
mask, to AnyText2. The result is a clean edit at the target's
natural width, without the gibberish-fill artifact common to
mask-based editors.

\subsection{Stage 4: Propagation}
\label{sec:method:s4}

S4 adapts the single edited ROI produced by S3 to match each
target frame's appearance. Two effects dominate: lighting and
blur.

\paragraph{Lighting correction.}
We adopt the Lighting Correction Module (LCM) structure from
STRIVE~\citep{subramanian2021strive}: inpaint the source text away
from both the reference and the target canonical ROIs (yielding
two clean backgrounds), compute a per-pixel multiplicative ratio
map between them, smooth with a Gaussian kernel, and multiply the
edited reference ROI by this ratio. Operations are carried out in
the log domain for numerical stability. Background isolation uses
one of two interchangeable inpainters (Section~\ref{sec:method:inpainter});
in our default configuration Hi-SAM~\citep{ye2024hisam} is used
here. When inpainted backgrounds are unavailable, S4 falls back to
a simpler global histogram match on the luminance channel.

\paragraph{Blur adaptation.}
A correctly lit edit can still look out of place if its sharpness
does not match the frame. Camera and object motion in video
produce frame-dependent blur that static compositing leaves
unmodeled. We address this with a Blur Prediction Network (BPN),
described in Section~\ref{sec:method:bpn}. Given the reference
canonical ROI and a small temporal window of target canonical
ROIs, the BPN predicts parameters of an anisotropic Gaussian blur
kernel and a residual weight, which are then applied to the
lighting-corrected edit via a differentiable blur operator. Each
frame thus receives an edit that is blur-matched to its local
temporal neighborhood.

S4 produces, for each frame, a \texttt{PropagatedROI} object
carrying the adapted pixels and a feathered alpha mask that is
approximately~$1$ in the interior and falls smoothly to~$0$ over
the outer $10\%$ of each dimension. The feathering ensures that
downstream compositing does not produce visible seams.

\subsection{Stage 5: Revert and Compositing}
\label{sec:method:s5}

S5 is the inverse of S2 followed by compositing. For each
detection we warp the corresponding \texttt{PropagatedROI} back to
the frame's perspective using $H_{\mathrm{from\_frontal}}$; in
practice we warp only to the bounding box of the target quad
(expanded by 5\%) rather than to the full frame, which is both
faster and more memory-efficient. Optional temporal smoothing can
be applied to the projected corner positions across a short
window, operating on corners rather than on raw $3\times 3$
matrices so that averaging remains geometrically meaningful.

Before compositing, S5 optionally runs a second pass of inpainting
over an expanded version of the target bounding box in the
original frame. This step cleans any residual source-text boundary
that falls outside the tight text mask of S4. The two inpainter
call sites---S4 background isolation and S5 pre-composite
cleanup---are independent instances with independent
configurations, allowing, for example, a Hi-SAM backend in S4
paired with an SRNet backend in S5.

Compositing itself is feathered alpha blending by default:
$F' = F\,(1-\alpha) + R\,\alpha$, where $F$ is the original frame,
$R$ is the warped propagated ROI, and $\alpha$ is the feathered
alpha mask from S4. A Poisson-blending
fallback~\citep{perez2003poisson} based on \texttt{cv2.seamlessClone}
is available for cases where gradient-domain seam matching is
preferred, but with proper feathering it is rarely needed.

\subsection{Blur Prediction Network}
\label{sec:method:bpn}

The BPN follows STRIVE's~\citep{subramanian2021strive} formulation:
the residual blur operator of Equation~\ref{eq:bpn} below, the
four-parameter anisotropic-Gaussian parameterization
$(\sigma_x, \sigma_y, \rho, w)$, the multi-target temporal batching,
and the two-stage (synthetic-supervised then self-supervised)
training schedule are all taken from STRIVE. Since STRIVE's code is
unreleased, every line of our BPN is new---but the algorithmic idea
is theirs, and we credit it as such. Our own contributions are
limited to three choices: (i)~the specific network architecture
(ResNet18 with tiled-weight channel-expansion initialization,
described below); (ii)~the training data preparation, which relies
on the Alignment Refiner's S2-refined homographies
(Section~\ref{sec:method:refiner}) and was generated through our
own pipeline; and (iii)~one inference-time robustness fix to the
blur operator for very small ROIs.

The BPN takes the concatenation of a reference canonical ROI and
$N$ target canonical ROIs along the channel dimension---a total of
$3(N+1)$ channels; we use $N=3$, giving $12$ channels at
$64\times 128$ input resolution---and predicts the blur parameters
$(\sigma_x, \sigma_y, \rho, w)$ for each of the $N$ targets. The
network uses a standard ResNet18~\citep{he2016resnet} backbone
whose first convolutional layer is replaced by a 12-channel input
convolution initialized by tiling the ImageNet-pretrained
three-channel weights across the expanded fan-in. The final
feature map is spatially averaged and passed through a two-layer
fully-connected head that emits $4N$ parameters; softplus and
$\tanh$ activations constrain $\sigma$ to be positive and
$(\rho, w)$ to lie in~$[-1, 1]$. The total parameter count is
approximately $1.24$M.

At inference, the predicted parameters parameterize a
differentiable blur operator:
\begin{equation}
  I_{\mathrm{out}} = (1 + w)\,I_{\mathrm{ref}} - w\,(I_{\mathrm{ref}} \ast G_{\sigma_x, \sigma_y, \rho}),
  \label{eq:bpn}
\end{equation}
where $G_{\sigma_x, \sigma_y, \rho}$ is an oriented anisotropic
Gaussian kernel. A small but important detail: we use replicate
padding rather than reflect padding in the blur convolution,
because reflect padding requires
$\mathrm{pad} < \mathrm{input\_size}$ in each dimension, which
fails on real-world tracks as small as 10--15~pixels tall when the
kernel is 41~pixels wide. Replicate has no such constraint and
gives equivalent halo suppression. This is an inference-time
robustness fix: training inputs are always $64\times 128$ and so
are unaffected.

Training proceeds in two stages. In Stage~1 (supervised), we
generate synthetic blur pairs by applying a known oriented
Gaussian blur with random $(\sigma_x, \sigma_y, \rho, w)$ to a
clean reference ROI and supervise on parameter MSE. This teaches
the network the parameter space without risking overfitting to
noisy real data. In Stage~2 (self-supervised), we draw two random
distinct frames from the same track---with no temporal-proximity
constraint, since in deployment the reference may be hundreds of
frames from the target---and supervise with a Charbonnier
reconstruction loss between the predicted-blurred reference and
the target, augmented by a temporal consistency term on adjacent
frame pairs.

Dataset preparation is a key detail. Pixel-level training requires
that reference and target canonical ROIs be geometrically aligned;
residual tracker error otherwise manifests as apparent blur, and
the network wastes capacity explaining geometry as blur. We
generate the training dataset through the pipeline's own
data-generation tool following the same S1$\rightarrow$S2 flow
described in Section~\ref{sec:method:refiner} for the Alignment
Refiner---PaddleOCR detection, CoTracker3 tracking, and
frontalization into per-track canonical rectangles---but here we
additionally compose the trained Alignment Refiner's $\Delta H$
into each frame's homography, so the canonical ROIs handed to the
BPN are already geometrically stabilized. This ensures that the
signal the network sees at training time matches the signal it
will see at inference time, with the refiner's contribution
absorbed as an alignment prior rather than an additional source of
pixel-level noise.

\subsection{Alignment Refiner}
\label{sec:method:refiner}

The Alignment Refiner takes a six-channel input (two RGB canonical
ROIs concatenated along the channel axis) at $64\times 128$
resolution and predicts eight real numbers: four $(dx, dy)$ corner
offsets. The network body is a deliberately small convolutional
encoder---four Conv-BN-ReLU blocks at strides of~$2$ with channel
widths $\{32, 64, 96, 128\}$, giving $16\times$ total
downsampling---followed by a flattened fully-connected head of
dimension $4096 \rightarrow 256 \rightarrow 8$. The final layer is
initialized with standard deviation $10^{-3}$ so that the
predicted residual is approximately zero at initialization. The
total parameter count is approximately $1.24$M, of which roughly
$85\%$ lie in the $4096 \rightarrow 256$ fully-connected layer.

The predicted corner offsets are converted into a $3\times 3$
residual homography $\Delta H$ via a differentiable direct linear
transform (DLT) solve, which allows gradients to flow through the
geometric parameterization during training.

\paragraph{Dataset generation.}
The real-pair half of the training data is generated by the
pipeline itself, which means the refiner is trained on exactly
the distribution of tracker error it will see at deployment rather
than on a synthetic proxy. Concretely, we run S1---PaddleOCR
detection followed by CoTracker3 online tracking---on a collection
of real videos, obtaining per-track quad trajectories. For each
track, we then apply S2's per-frame frontalization homography to
warp the quad at every frame into the track's canonical rectified
space, yielding one canonical ROI per frame. Because every ROI in
a given track sits in the same canonical coordinate system, any
two of them \emph{should} be pixel-aligned; in practice they
differ by the residual CoTracker and homography-fit error that the
refiner is meant to correct. Sampling two random frames from the
same track therefore produces a training pair whose disagreement
is exactly the signal we want to supervise on. The preprocessing
is run once offline and stored as a flat directory of per-track
ROIs; the training loader sees a pool of tracks and samples pairs
on the fly. For Type~A (synthetic) samples we reuse the same
directory but draw a single ROI and apply a synthetic corner
perturbation instead of picking a second real frame.

We train with a mixed-batch, two-stage schedule. The dataset
exposes two sample types. \emph{Type~A} is synthetic: we take a
clean ROI, apply a random $\pm 8$-pixel corner perturbation, and
supervise the network to predict the inverse perturbation. The
loss is masked Charbonnier on RGB plus a small corner-L2
regularizer. \emph{Type~B} is real: two random distinct frames
from the same track, with no temporal-proximity constraint (a
deployment reference may be hundreds of frames from its target,
and the network must not rely on small-motion priors). The
loss is a masked normalized cross-correlation (NCC) on luminance
(shift- and scale-invariant) plus a masked Charbonnier on Sobel
magnitude (structure-preserving), with the same corner
regularizer. A per-item \texttt{real\_pair\_fraction} selects
between the two sample types. Stage~1 trains with 100\%
Type~A until the synthetic error plateaus; Stage~2 resumes from
that checkpoint with 70\% Type~B / 30\% Type~A.

At inference, a sanity gate rejects any prediction whose maximum
corner offset exceeds 16~pixels; rejected detections silently fall
back to the unrefined baseline homography. We additionally log at
\texttt{INFO} when at least $10\%$ of a track's predictions
trigger the rejection gate, which functions as an early warning
that the reference frame is an outlier for that track.

\subsection{Inpainter backends}
\label{sec:method:inpainter}

Three of the pipeline's call sites require removing existing text
from an ROI and filling the background plausibly: S3's
adaptive-mask pre-inpaint, S4's background isolation for LCM, and
S5's pre-composite edge cleanup. We support two interchangeable
backends, summarized in Table~\ref{tab:inpainter}:
SRNet~\citep{wu2019srnet} reused as a learned text-removal
network, and a segmentation-based pipeline that combines
Hi-SAM~\citep{ye2024hisam} for stroke-level text segmentation with
OpenCV's Navier--Stokes inpainting for the pixel fill. In our
default configuration Hi-SAM is used in S4, with a small
three-pixel mask dilation to absorb anti-aliased stroke edges; its
remaining failure mode on blurred text is documented in
Table~\ref{tab:inpainter}. Each call site instantiates its
inpainter independently, so the three can use different backends
or different configurations without affecting one another.

\begin{table}[t]
  \centering
  \caption{Inpainter backend tradeoffs.}
  \label{tab:inpainter}
  \begin{tabular}{lp{0.28\linewidth}p{0.38\linewidth}}
    \toprule
    & SRNet & Hi-SAM + NS inpaint \\
    \midrule
    Peak VRAM at inference  & $\sim$200--300\,MB
        & $\sim$5--7\,GB in fp32, dominated by SAM ViT-L's
          $1024\times 1024$ forward pass (\textasciitilde{}1.2\,GB of
          weights plus several GB of attention activations across
          24 transformer blocks). Drops roughly $2\times$ in fp16 \\
    Speed                   & Faster (small CNN forward)
        & Slower (SAM ViT-L backbone dominates per-ROI cost) \\
    Style generalization    & Training-distribution only
        & Generic (SAM-based) \\
    Typical failure         & Residual faint strokes where the text
        removal is incomplete
        & Heavily blurred or low-contrast text is under-segmented,
          leaving a faint, thickened ghost of the original text in
          the inpainted background \\
    Mitigation              & ---
        & 3\,px mask dilation (absorbs anti-aliased stroke edges;
          does not recover text the segmenter missed entirely) \\
    \bottomrule
  \end{tabular}
\end{table}

\subsection{Method summary}
\label{sec:method:summary}

The methodology above rests on three design choices that together
produce the pipeline's core properties.

\begin{itemize}[leftmargin=*]
  \item \textbf{Explicit geometric normalization.} Frontalizing
    every detection into a shared canonical rectangle decouples
    text editing from perspective distortion, and the learned
    Alignment Refiner absorbs the small remaining misalignment
    that optical-flow trackers leave. All downstream stages
    therefore operate in a pixel-aligned canonical space.
  \item \textbf{Stage-wise decomposition.} Text generation happens
    once per track on a single high-quality reference frame;
    propagation to every other frame is an
    appearance-adaptation problem (lighting and blur) rather than
    a re-editing problem. This avoids the per-frame inference
    cost of generative video editors and removes a major source
    of temporal flicker.
  \item \textbf{Appearance-aware propagation.} The Lighting
    Correction Module handles spatially-varying illumination via
    per-pixel ratio maps computed on inpainted backgrounds;
    the Blur Prediction Network handles per-frame motion and
    focus blur via a compact learned predictor. The two target
    complementary axes of appearance variation.
\end{itemize}

Together, this decomposition produces temporally consistent
cross-language scene text replacement without per-frame generative
inference, while keeping every intermediate representation
(canonical ROIs, homography matrices, lighting ratio maps, blur
parameters) inspectable and individually controllable.
